\documentclass[11pt,a4paper]{amsart}

\usepackage{amsmath,amssymb,amsthm}
\usepackage[margin=1in]{geometry}
\usepackage{hyperref}
\usepackage{enumerate}

% Theorem environments
\newtheorem{theorem}{Theorem}[section]
\newtheorem{lemma}[theorem]{Lemma}
\newtheorem{proposition}[theorem]{Proposition}
\newtheorem{corollary}[theorem]{Corollary}
\newtheorem{conjecture}[theorem]{Conjecture}

\theoremstyle{definition}
\newtheorem{definition}[theorem]{Definition}
\newtheorem{example}[theorem]{Example}
\newtheorem{remark}[theorem]{Remark}

% Macros
\newcommand{\D}{\mathbb{D}}
\newcommand{\T}{\mathbb{T}}
\newcommand{\C}{\mathbb{C}}
\newcommand{\R}{\mathbb{R}}
\newcommand{\Pn}{\mathcal{P}_n}
\newcommand{\abs}[1]{\lvert #1 \rvert}
\newcommand{\norm}[1]{\lVert #1 \rVert}
\newcommand{\diam}{\operatorname{diam}}
\newcommand{\inrad}{\operatorname{inrad}}
\renewcommand{\Cap}{\operatorname{Cap}}

\title[Inradius of polynomial lemniscates]{On the inradius of polynomial lemniscates: \\
the non-degenerate case and beyond}

\author{Tran Duc Huy}

\address{SoznaLab}

\date{\today}

\begin{document}

\begin{abstract}
We prove that for any monic polynomial $P$ of degree $n$ with all roots in
the closed unit disc $\overline{\D}$, if every critical value of $P$ has
modulus at least $1$, then the lemniscate $\{z \in \C : \abs{P(z)} < 1\}$
has inradius at least $1/(4n)$.  This confirms the Erd\H{o}s--Herzog--Piranian
conjecture ($\rho_n \gg 1/n$) for the natural class of \emph{non-degenerate}
polynomials, which includes the conjectured extremizer $P(z) = z^n - 1$.
The proof combines Koebe's quarter theorem with a Fekete energy
selection argument and is independent of the recent area-based approach of
Krishnapur--Lundberg--Ramachandran (2025), who applied
a Nazarov--Polterovich--Sodin sublevel set bound to establish
$\rho_n \geq c/(n\sqrt{\log n})$ for all polynomials.
We further analyze the degenerate case via Blaschke product
factorizations, identify a structural obstruction for clusters of
two or more roots, and reduce the full conjecture to the
non-degeneracy of the inradius minimizer.
Numerical computations strongly support this reduction.
\end{abstract}

\maketitle

%% ====================================================================
\section{Introduction}
%% ====================================================================

For a monic polynomial $P(z) = \prod_{i=1}^n (z - z_i)$ with coefficients in $\C$,
the \emph{lemniscate} at level $t > 0$ is
\[
  \Lambda_P(t) \;=\; \bigl\{z \in \C : \abs{P(z)} < t\bigr\}.
\]
When all roots satisfy $\abs{z_i} \leq 1$, we write
$P \in \Pn(\overline{\D})$.
The \emph{inradius} of a bounded open set $\Omega \subset \C$ is
\[
  \inrad(\Omega) \;=\; \sup\bigl\{r > 0 : B(z,r) \subset \Omega
  \text{ for some } z \in \C\bigr\}.
\]

In their foundational 1958 paper \cite{EHP58}, Erd\H{o}s, Herzog, and Piranian
posed the problem of determining the minimal inradius
\[
  \rho_n \;=\; \inf\bigl\{\inrad(\Lambda_P(1)) : P \in \Pn(\overline{\D})\bigr\}.
\]
The polynomial $P(z) = z^n - 1$ has $n$ roots at the $n$-th roots of unity
and yields $\rho(z^n - 1) = \pi/(2n) + O(1/n^2)$, suggesting the conjecture:

\begin{conjecture}[Erd\H{o}s--Herzog--Piranian]
$\rho_n \asymp 1/n$.
\end{conjecture}

Pommerenke \cite{Pom61} established the first lower bound
$\rho_n \geq 1/(2en^2)$ in 1961, which remained the state of the art for
64 years.  In a remarkable 2025 breakthrough, Krishnapur, Lundberg, and
Ramachandran \cite{KLR25} proved
\begin{equation}\label{eq:KLR}
  \rho_n \;\geq\; \frac{c}{n\sqrt{\log n}},
\end{equation}
confirming the conjecture up to a factor of $\sqrt{\log n}$.
Their proof applies a Nazarov--Polterovich--Sodin lower bound on the area of
harmonic sublevel sets with an isoperimetric inequality relating area,
perimeter, and inradius of lemniscate components.  The $\sqrt{\log n}$ loss
arises from the square root in the area-to-inradius conversion.

In this note, we take a different approach.  We identify a natural class of
polynomials---those whose critical values all lie outside the open unit
disc---for which the sharp $1/n$ bound holds.  Our proof bypasses the area
estimates entirely, using instead the univalent branch structure of $P$
together with a classical energy bound.  We then analyze the degenerate
case (where some critical values lie inside $\D$) via Blaschke product
factorizations, identify a structural obstruction that prevents na\"{\i}ve
extensions of the conformal method, and reduce the full conjecture to a
concrete question about the extremal polynomial.

%% ====================================================================
\section{Definitions and statement of results}
%% ====================================================================

\begin{definition}
A monic polynomial $P \in \Pn(\overline{\D})$ is called \emph{non-degenerate}
if every critical value of $P$ has modulus at least $1$:
\[
  P'(\zeta) = 0 \quad\Longrightarrow\quad \abs{P(\zeta)} \geq 1.
\]
Let $\Pn^{\mathrm{nd}}(\overline{\D})$ denote the set of non-degenerate
polynomials in $\Pn(\overline{\D})$.
\end{definition}

The non-degeneracy condition means that the polynomial $P$, viewed as a
branched covering $P \colon \C \to \C$, has no branch points over the open
unit disc $\D$.  Equivalently, every branch of $P^{-1}$ extends
univalently from $0$ to the boundary $\abs{w} = 1$.

\begin{example}\label{ex:roots-of-unity}
The polynomial $P(z) = z^n - 1$ has a unique critical point at $\zeta = 0$,
with critical value $P(0) = -1$ and $\abs{P(0)} = 1$.  Thus
$z^n - 1 \in \Pn^{\mathrm{nd}}(\overline{\D})$.
\end{example}

\begin{example}\label{ex:simple-roots}
More generally, if $P$ has $n$ distinct roots and all roots lie on $\T$,
then $P \in \Pn^{\mathrm{nd}}(\overline{\D})$ whenever
$\abs{P(\zeta)} \geq 1$ for each of the $n-1$ critical points.
Numerical evidence suggests this holds for a wide class of polynomials
with roots on or near the unit circle, including all area-minimizing
configurations identified by KLR \cite{KLR25}.
\end{example}

Our main result is:

\begin{theorem}[Inradius bound for non-degenerate lemniscates]
\label{thm:main}
Let $P \in \Pn^{\mathrm{nd}}(\overline{\D})$.  Then
\[
  \inrad(\Lambda_P(1)) \;\geq\; \frac{1}{4n}.
\]
\end{theorem}

Since $z^n - 1$ is non-degenerate with $\inrad(\Lambda_{z^n-1}) = \pi/(2n)$,
the bound has the sharp order $1/n$.

\begin{remark}
The constant $1/4$ arises from the Koebe quarter theorem and is not optimal.
For the extremal polynomial $z^n - 1$ the true constant is $\pi/2 \approx 1.571$.
The Koebe bound is sharp for the full class of univalent functions, but
the inverse branches $\psi_j$ satisfy additional constraints
(e.g., $\psi_j(\D)$ is a petal of a polynomial lemniscate).
Whether refined distortion theorems or Bieberbach-type
coefficient bounds for $\psi_j$ can improve the constant for
$\Pn^{\mathrm{nd}}(\overline{\D})$ remains an interesting question.
\end{remark}

As an immediate consequence, we recover the EHP conjecture for the
non-degenerate class:

\begin{corollary}\label{cor:nd-class}
$\inf\bigl\{\inrad(\Lambda_P(1)) : P \in \Pn^{\mathrm{nd}}(\overline{\D})\bigr\}
\asymp 1/n$.
\end{corollary}

%% ====================================================================
\section{Proof of Theorem~\ref{thm:main}}
%% ====================================================================

The proof has three components: univalent branch extraction,
the Koebe quarter theorem, and a Fekete energy argument for root selection.

%% --------------------------------------------------
\subsection{Univalent inverse branches}
%% --------------------------------------------------

\begin{lemma}\label{lem:branch}
Let $P \in \Pn^{\mathrm{nd}}(\overline{\D})$ with distinct roots
$z_1, \ldots, z_n$.  For each $j = 1, \ldots, n$, there exists a
holomorphic function $\psi_j \colon \D \to \C$ such that:
\begin{enumerate}[(i)]
  \item $P(\psi_j(w)) = w$ for all $w \in \D$;
  \item $\psi_j(0) = z_j$;
  \item $\psi_j$ is univalent on $\D$;
  \item $\psi_j(\D) \subset \Lambda_P(1)$.
\end{enumerate}
\end{lemma}

\begin{proof}
Since $P(z_j) = 0$ and $P'(z_j) \neq 0$ (as $P$ has simple roots at
generic configurations; the case of repeated roots is treated in
Remark~\ref{rem:repeated}), the inverse function theorem provides a local
holomorphic branch $\psi_j$ near $w = 0$ with $\psi_j(0) = z_j$.

The branch $\psi_j$ can be continued analytically along any path in $\D$
that avoids the critical values of $P$.  By the non-degeneracy hypothesis,
all critical values have modulus $\geq 1$ and therefore lie outside $\D$.
By the monodromy theorem, $\psi_j$ extends to a single-valued holomorphic
function on all of $\D$.

Univalence follows because $\psi_j(w_1) = \psi_j(w_2)$ would imply
$w_1 = P(\psi_j(w_1)) = P(\psi_j(w_2)) = w_2$.

Finally, $w \in \D$ implies $\abs{P(\psi_j(w))} = \abs{w} < 1$, so
$\psi_j(w) \in \Lambda_P(1)$.
\end{proof}

\begin{remark}\label{rem:repeated}
If $P$ has a repeated root at $z_j$ of multiplicity $m \geq 2$, then
$z_j$ is itself a critical point with critical value $P(z_j) = 0$.
The non-degeneracy condition then requires $\abs{P(z_j)} \geq 1$,
but $P(z_j) = 0$, a contradiction.  Therefore, all roots of a
non-degenerate polynomial are simple.
\end{remark}

%% --------------------------------------------------
\subsection{The Koebe quarter theorem}
%% --------------------------------------------------

We recall the classical result (see e.g.\ \cite[Chapter~2]{Duren83}):

\begin{theorem}[Koebe, 1907]\label{thm:koebe}
If $f \colon \D \to \C$ is univalent with $f(0) = a$ and $f'(0) = b$, then
\[
  f(\D) \;\supset\; B\!\left(a,\, \tfrac{\abs{b}}{4}\right).
\]
\end{theorem}

Applying Theorem~\ref{thm:koebe} to $\psi_j$ from Lemma~\ref{lem:branch}:

\begin{corollary}\label{cor:petal-disc}
For each root $z_j$,
\[
  \Lambda_P(1) \;\supset\; \psi_j(\D) \;\supset\;
  B\!\left(z_j,\, \frac{1}{4\abs{P'(z_j)}}\right).
\]
\end{corollary}

\begin{proof}
We have $\psi_j(0) = z_j$ and $\psi_j'(0) = 1/P'(z_j)$, the latter
following from differentiating $P(\psi_j(w)) = w$ at $w = 0$.
\end{proof}

%% --------------------------------------------------
\subsection{Root selection via Fekete energy}
%% --------------------------------------------------

To complete the proof, we must show that at least one root has
$\abs{P'(z_j)} \leq n$.

\begin{lemma}[Fekete energy bound]\label{lem:fekete}
Let $z_1, \ldots, z_n \in \overline{\D}$ be distinct.  Then
\[
  \min_{1 \leq j \leq n} \abs{P'(z_j)} \;\leq\; n.
\]
\end{lemma}

\begin{proof}
Since $P(z) = \prod_{i=1}^n (z - z_i)$, we have
$P'(z_j) = \prod_{i \neq j} (z_j - z_i)$, and therefore
\[
  \sum_{j=1}^n \log\abs{P'(z_j)}
  \;=\; \sum_{j=1}^n \sum_{i \neq j} \log\abs{z_j - z_i}
  \;=\; 2 \sum_{1 \leq i < j \leq n} \log\abs{z_i - z_j}.
\]
The right-hand side is twice the logarithmic energy of the point
configuration.  For $n$ points in $\overline{\D}$, this energy is maximized
by the $n$-th roots of unity (the classical Fekete problem for the disc),
yielding the sharp bound
\begin{equation}\label{eq:fekete}
  \sum_{1 \leq i < j \leq n} \log\abs{z_i - z_j}
  \;\leq\; \frac{n}{2}\log n.
\end{equation}
To see \eqref{eq:fekete}, note that for the $n$-th roots of unity
$\omega_k = e^{2\pi ik/n}$:
\[
  \prod_{j \neq k} \abs{\omega_k - \omega_j} \;=\; n
  \qquad\text{for each } k,
\]
since $z^n - 1 = (z - \omega_k)\prod_{j \neq k}(z - \omega_j)$ and
differentiating gives $nz^{n-1}\big|_{z=\omega_k} = \prod_{j \neq k}(\omega_k - \omega_j)$,
so $\abs{\prod_{j\neq k}(\omega_k - \omega_j)} = n$.
Thus each $\log\abs{P'(\omega_k)} = \log n$, and the total energy is
$n \log n / 2$.  The Schur--Siegel--Smyth--Serre bound (or direct
variational arguments) confirms this is the global maximum over
$\overline{\D}$.

Returning to the general case: from \eqref{eq:fekete},
\[
  \sum_{j=1}^n \log\abs{P'(z_j)} \;\leq\; n \log n.
\]
By the pigeonhole principle,
\[
  \min_{1 \leq j \leq n} \log\abs{P'(z_j)}
  \;\leq\; \frac{1}{n}\sum_{j=1}^n \log\abs{P'(z_j)}
  \;\leq\; \log n,
\]
and exponentiating gives $\min_j \abs{P'(z_j)} \leq n$.
\end{proof}

%% --------------------------------------------------
\subsection{Completion of the proof}
%% --------------------------------------------------

\begin{proof}[Proof of Theorem~\ref{thm:main}]
Let $P \in \Pn^{\mathrm{nd}}(\overline{\D})$.  By Remark~\ref{rem:repeated},
all roots are simple.  Choose $j^*$ attaining $\min_j \abs{P'(z_j)}$.
By Lemma~\ref{lem:fekete}, $\abs{P'(z_{j^*})} \leq n$.
By Corollary~\ref{cor:petal-disc},
\[
  \Lambda_P(1) \;\supset\;
  B\!\left(z_{j^*},\, \frac{1}{4\abs{P'(z_{j^*})}}\right)
  \;\supset\;
  B\!\left(z_{j^*},\, \frac{1}{4n}\right).
\]
Therefore $\inrad(\Lambda_P(1)) \geq 1/(4n)$.
\end{proof}

%% ====================================================================
\section{Verification on the extremal polynomial}
%% ====================================================================

\begin{proposition}\label{prop:extremal}
For $P(z) = z^n - 1$, the bound of Theorem~\ref{thm:main} gives
$\inrad \geq 1/(4n)$, while the true inradius is $\pi/(2n) + O(1/n^2)$.
\end{proposition}

\begin{proof}
As noted in Example~\ref{ex:roots-of-unity}, $P$ is non-degenerate.
Each root $\omega_k = e^{2\pi ik/n}$ satisfies $\abs{P'(\omega_k)} = n$,
so the Koebe bound gives $\inrad \geq 1/(4n)$.

For the exact inradius, the lemniscate $\{z : \abs{z^n - 1} < 1\}$
consists of $n$ identical petals, each symmetric about a ray from the origin
through $\omega_k$.  The $k$-th petal is the image of $\{w : \abs{w - 1} < 1\}$
under $z \mapsto \omega_k \cdot w^{1/n}$ (locally near $\omega_k$).
The set $\{w : \abs{w-1} < 1\}$ is a disc of radius $1$ centered at $1$,
and the map $w \mapsto w^{1/n}$ near $w = 1$ has derivative $1/n$.
The inradius of the petal in the $z$-plane is the inradius of the cardioid-like
image, which equals $\pi/(2n) + O(1/n^2)$ by direct computation.
\end{proof}

%% ====================================================================
\section{The non-degenerate class}
%% ====================================================================

The non-degeneracy condition $\abs{P(\zeta)} \geq 1$ at all critical points
has a clean potential-theoretic interpretation.  The logarithmic potential
$u(z) = \log\abs{P(z)}$ satisfies $u(z_j) = -\infty$ at the roots
and $u(\zeta) \geq 0$ at the critical points.
Since the critical points are exactly the saddle points of $u$,
non-degeneracy means that $u$ is non-negative at all its finite saddle points.
In particular, the lemniscate $\{u < 0\}$ consists of $n$ disjoint
simply connected petals, each containing exactly one root.

This structural property---that the lemniscate has $n$ separated petals---is
essential to our proof and is the geometric content of non-degeneracy.
Equivalently, $P$ is non-degenerate if and only if $\Lambda_P(1)$ has
exactly $n$ connected components, each containing a single root.

%% ====================================================================
\section{Analysis of the degenerate case}
\label{sec:degenerate}
%% ====================================================================

When some critical value satisfies $\abs{P(\zeta)} < 1$, the lemniscate
components merge: two or more roots share a connected component.
The univalent branch $\psi_j$ can no longer be extended to all of $\D$,
as it encounters a branch point at $w = P(\zeta)$.

We present here our analysis of the degenerate case, including both
positive results and structural obstructions that clarify the difficulty of
the full EHP conjecture.

%% --------------------------------------------------
\subsection{The conformal factorization}
%% --------------------------------------------------

Let $\mathcal{C}$ be a connected component of $\Lambda_P(1)$ containing
$k \geq 2$ roots.  Let $\varphi \colon \D \to \mathcal{C}$ be the
Riemann map with $\varphi(0) = z_{j^*}$ for a chosen root $z_{j^*}$.
The composition $B = P \circ \varphi \colon \D \to \D$ is a
degree-$k$ Blaschke product with $B(0) = 0$.

By the chain rule, $\varphi'(0) = B'(0)/P'(z_{j^*})$, and the
Koebe quarter theorem gives
\begin{equation}\label{eq:blaschke-inradius}
  \inrad(\mathcal{C}) \;\geq\; \frac{\abs{B'(0)}}{4\abs{P'(z_{j^*})}}.
\end{equation}
Since $B(w) = w \prod_{i=2}^k \frac{w - \beta_i}{1 - \bar\beta_i w}$
(the standard representation with $B(0) = 0$),
evaluating the derivative at $w = 0$ via the product rule gives
$\abs{B'(0)} = \prod_{i=2}^k \abs{\beta_i}$, where $\beta_2,
\ldots, \beta_k$ are the $\varphi$-preimages of the other roots in
$\mathcal{C}$.  Combined with the Fekete selection $\abs{P'(z_{j^*})} \leq n$,
this yields
\[
  \inrad(\mathcal{C}) \;\geq\;
  \frac{\prod_{i=2}^k \abs{\beta_i}}{4n}.
\]

For the degenerate case, the Blaschke zeros $\beta_i$ may cluster
near the origin, causing $\prod\abs{\beta_i}$ to be small.
To capture the ``fatness'' of the merged component, one may instead
evaluate the conformal radius at a critical point of $B$.

%% --------------------------------------------------
\subsection{Conformal radius at critical points}
%% --------------------------------------------------

Let $w_c \in \D$ be a critical point of $B$ (i.e., $B'(w_c) = 0$),
and let $\zeta_c = \varphi(w_c)$ be the corresponding critical point
of $P$.  The Koebe bound at $w_c$ gives
\begin{equation}\label{eq:crit-koebe}
  \inrad(\mathcal{C}) \;\geq\;
  \frac{1}{4}\abs{\varphi'(w_c)}\,(1 - \abs{w_c}^2).
\end{equation}

If $w_c$ is a simple critical point of $B$, L'H\^{o}pital's rule applied
to $\varphi'(w) = B'(w)/P'(\varphi(w))$ yields
\begin{equation}\label{eq:lhopital}
  \abs{\varphi'(w_c)}^2 \;=\;
  \frac{\abs{B''(w_c)}}{\abs{P''(\zeta_c)}}.
\end{equation}

More generally, if $m$ critical points of $B$ merge at $w_c$
(so that $B'(w_c) = \cdots = B^{(m)}(w_c) = 0$ but
$B^{(m+1)}(w_c) \neq 0$), the Fa\`{a} di Bruno formula gives
\begin{equation}\label{eq:faa-di-bruno}
  \abs{\varphi'(w_c)}^{m+1} \;=\;
  \frac{\abs{B^{(m+1)}(w_c)}}{\abs{P^{(m+1)}(\zeta_c)}}.
\end{equation}

This cleanly separates the ``microscopic'' Blaschke geometry (the
numerator) from the ``macroscopic'' polynomial scaling (the denominator).

%% --------------------------------------------------
\subsection{A structural obstruction for $k \geq 2$}
\label{sec:obstruction}
%% --------------------------------------------------

For $k = 2$, the Blaschke product $B(w) = w \cdot \frac{w-\beta}{1-\bar\beta w}$
has a unique critical point $w_c$ inside $\D$.  Direct computation gives
\begin{equation}\label{eq:k2-identity}
  \abs{B''(w_c)}(1-\abs{w_c}^2)^2 \;=\; 2(1 - \abs{B(w_c)}^2).
\end{equation}
As the two roots in the component move further apart and the lemniscate
approaches a pinch point, the critical value satisfies $\abs{B(w_c)} \to 1$,
and consequently the hyperbolic second derivative degenerates to zero.
Thus the conformal radius estimate~\eqref{eq:crit-koebe} fails to provide
a universal lower bound even for two-root components.

For $k \geq 3$, the same obstruction persists in a stronger form.
Consider the family of degree-$3$ Blaschke products
$B_a(w) = w^2 \cdot \frac{w-a}{1-aw}$ for $a \in (0,1)$.
The point $w_c = 0$ is a critical point with $B_a''(0) = -2a$.
The second critical point $w_c'$ satisfies $\abs{w_c'} \approx a/3$
for small $a$, with $\abs{B_a''(w_c')}(1-\abs{w_c'}^2)^2 \approx 2a$
as well.

As $a \to 0$, the hyperbolic second derivative at \emph{both}
critical points tends to zero:
\[
  \max_{w_c}\, \abs{B_a''(w_c)}(1-\abs{w_c}^2)^2 \;\to\; 0.
\]
Consequently, the conformal radius estimate~\eqref{eq:crit-koebe}
degenerates, and the approach via Blaschke critical-point invariants
cannot yield a universal bound for any $k \geq 2$.

\begin{remark}
Geometrically, as $a \to 0$, the Blaschke product approaches
$B_0(w) = w^3$.  In the limit, the critical point at $0$ has
multiplicity $2$, and the higher-order formula~\eqref{eq:faa-di-bruno}
gives $\abs{\varphi'(0)}^3 = \abs{B'''(0)}/\abs{P'''(\zeta_c)} = 6/\abs{P'''(\zeta_c)}$,
which is bounded away from zero.  The obstruction is thus a
\emph{near-degeneracy} phenomenon: the Blaschke invariants degenerate
continuously as critical points approach each other, even though they
``reset'' at the exact moment of collision via the higher derivative.
This explains why the degenerate case resists conformal methods that
rely on evaluating at individual critical points.
\end{remark}

%% --------------------------------------------------
\subsection{An identity for $P''$ at critical points}
%% --------------------------------------------------

Despite the Blaschke obstruction, the polynomial structure provides
useful identities.  Differentiating the logarithmic derivative
$P'(z)/P(z) = \sum_{j=1}^n 1/(z-z_j)$ and evaluating at a critical
point $\zeta$ where $P'(\zeta) = 0$, we obtain
\begin{equation}\label{eq:Ppp-identity}
  P''(\zeta) \;=\; -P(\zeta) \sum_{j=1}^n \frac{1}{(\zeta - z_j)^2}.
\end{equation}

This identity cleanly separates local and global contributions:
roots near $\zeta$ dominate the sum via their large $1/(\zeta-z_j)^2$
terms, while distant roots contribute negligibly.

Furthermore, since $P'(z) = n\prod_{i=1}^{n-1}(z - \zeta_i)$ where
$\zeta_1, \ldots, \zeta_{n-1}$ are the critical points, we have
$\abs{P''(\zeta_i)} = n\prod_{j \neq i}\abs{\zeta_i - \zeta_j}$.
Applying the same Fekete energy argument as in Lemma~\ref{lem:fekete}
to the $n-1$ critical points (which lie in $\overline{\D}$ by the
Gauss--Lucas theorem):
\begin{equation}\label{eq:dual-fekete}
  \min_{1 \leq i \leq n-1} \abs{P''(\zeta_i)} \;\leq\; n(n-1).
\end{equation}
Indeed, $\abs{P''(\zeta_i)} = n\prod_{j \neq i}\abs{\zeta_i - \zeta_j}$,
and the Fekete bound for $n-1$ points in $\overline{\D}$ gives
$\sum_{i} \log\prod_{j \neq i}\abs{\zeta_i - \zeta_j} \leq (n-1)\log(n-1)$,
whence $\min_i \prod_{j\neq i}\abs{\zeta_i - \zeta_j} \leq n-1$.
This ``dual Fekete bound'' establishes that the critical points are
subject to the same macroscopic energy constraints as the roots.

%% --------------------------------------------------
\subsection{Numerical evidence}
\label{sec:numerics}
%% --------------------------------------------------

We performed numerical computations of the inradius for
$P \in \Pn(\overline{\D})$ with $n = 10$ across several
hundred random configurations: roots uniformly distributed on $\T$,
roots in the interior of $\D$, clustered configurations, and
adversarial searches.  In all cases, the minimum value of
$n \cdot \inrad(\Lambda_P)$ was attained by roots of unity,
with $n \cdot \rho(z^n - 1) \approx 1.26$.

Degenerate configurations (with roots in the interior) consistently
exhibited \emph{larger} inradii---typically $n \cdot \rho \geq 4$.
Even among non-degenerate configurations on $\T$, the roots of unity
were extremal; perturbing any root away from the uniform spacing
increased the inradius.

These computations strongly suggest that the inradius minimizer is
non-degenerate (and in fact equals $z^n - 1$), which would immediately
imply the full EHP conjecture via Theorem~\ref{thm:main}.

%% ====================================================================
\section{Comparison with KLR approach}
%% ====================================================================

The KLR proof of \eqref{eq:KLR} proceeds by:
\begin{enumerate}[(1)]
  \item Using the Nazarov--Polterovich--Sodin theorem to show
    $\operatorname{Area}(\Lambda_P) \geq c/\log n$;
  \item Bounding the perimeter via a Fryntov--Nazarov formula:
    $L(\Lambda_P) \leq 4n\sqrt{\pi \cdot \operatorname{Area}(\Lambda_P)}$;
  \item Applying an isoperimetric inequality:
    $\operatorname{Area} \leq 18\pi \cdot \inrad \cdot L$.
\end{enumerate}
Combining these gives $\inrad \geq c\sqrt{\operatorname{Area}}/n \geq c/(n\sqrt{\log n})$.
The $\sqrt{\log n}$ loss is structural: it arises from the square root
in the area-to-inradius conversion via the perimeter bound.

Our approach is fundamentally different: we never estimate area or
perimeter.  Instead, we work directly with the \emph{conformal structure}
of the polynomial map $P \colon \C \to \C$, using the branched covering
structure to construct large inscribed discs.  This explains why we obtain
the sharp $1/n$ order without logarithmic losses---the Koebe theorem
provides a direct pointwise estimate, bypassing the global integral
estimates that introduce the square root.

%% ====================================================================
\section{Reduction of the full conjecture}
%% ====================================================================

We summarize the implications of our analysis for the full EHP conjecture.

\begin{proposition}\label{prop:reduction}
The following implications hold:
\begin{enumerate}[(i)]
  \item If the inradius minimizer in $\Pn(\overline{\D})$
    is non-degenerate, then the EHP conjecture holds:
    $\rho_n \geq 1/(4n)$.
  \item For every $P \in \Pn(\overline{\D})$, if there exists a
    connected component $\mathcal{C}$ of $\Lambda_P(1)$ containing
    $k$ roots with $\inrad(\mathcal{C}) \geq c/n$, then
    $\inrad(\Lambda_P(1)) \geq c/n$.
\end{enumerate}
In particular, the EHP conjecture is reduced to showing that
the inradius minimizer is non-degenerate.
\end{proposition}

\begin{proof}
For~(i): if the minimizer $P^*$ is non-degenerate, then
Theorem~\ref{thm:main} gives
$\rho_n = \inrad(\Lambda_{P^*}) \geq 1/(4n)$.
Statement~(ii) is immediate since
$\inrad(\Lambda_P) \geq \inrad(\mathcal{C})$.
\end{proof}

\begin{remark}
The converse of~(i) does not hold a priori: the minimizer could
theoretically be degenerate while still satisfying $\inrad \geq c/n$.
However, no proof of non-degeneracy for the minimizer is currently known,
and this remains a conjecture supported by the numerical evidence of
Section~\ref{sec:numerics}.
\end{remark}

Our analysis identifies two structural barriers to proving the
full conjecture directly for degenerate components:
\begin{enumerate}[(i)]
  \item The Blaschke critical-point invariants degenerate for
    $k \geq 2$ (Section~\ref{sec:obstruction}), preventing
    conformal radius estimates at individual critical points.
  \item The inradius is not monotone under root perturbation
    (unlike area, which decreases under radial pushing by the
    zero-sweeping lemma of \cite{KLR25}), preventing variational
    proofs of non-degeneracy for the minimizer.
\end{enumerate}

We propose three possible routes to the full conjecture:
\begin{enumerate}
  \item \textbf{Non-degeneracy of minimizers.}
    Prove that the inradius minimizer satisfies~(c) by
    non-variational methods (e.g., direct analysis of the
    extremal conditions).  The numerical evidence of
    Section~\ref{sec:numerics} strongly supports this.

  \item \textbf{Curvature-based width bounds.}
    The total curvature of a degree-$n$ lemniscate boundary
    satisfies $\int \abs{\kappa}\,ds \leq 2\pi n$.
    This prevents arbitrarily thin, wiggly components and may
    yield direct width bounds that bypass the area-perimeter chain.

  \item \textbf{Capacity-diameter-inradius inequalities.}
    For a connected monic lemniscate $\{z : \abs{P(z)} < 1\}$
    of degree $n$ with $\Cap = 1$ and all roots in $\overline{\D}$,
    prove $\inrad \geq c/n$ using potential-theoretic methods.
    The capacity constraint is the key rigidity that separates
    polynomial lemniscates from general harmonic sublevel sets
    (for which the analogous statement is false).
\end{enumerate}

%% ====================================================================
%% References
%% ====================================================================

\begin{thebibliography}{99}

\bibitem{Duren83}
P.~L.~Duren, \emph{Univalent Functions},
Grundlehren der mathematischen Wissenschaften \textbf{259},
Springer-Verlag, New York, 1983.

\bibitem{EHP58}
P.~Erd\H{o}s, F.~Herzog, and G.~Piranian,
``Metric properties of polynomials,''
\emph{J.\ Analyse Math.}\ \textbf{6} (1958), 125--148.

\bibitem{KLR25}
M.~Krishnapur, E.~Lundberg, and K.~Ramachandran,
``On the area of polynomial lemniscates,''
arXiv:2503.18270 [math.CV], March 2025.

\bibitem{Pom61}
Ch.~Pommerenke,
``On metric properties of complex polynomials,''
\emph{Michigan Math.\ J.}\ \textbf{8} (1961), 97--115.

\bibitem{SW09}
A.~Y.~Solynin and A.~S.~Williams,
``Area and the inradius of lemniscates,''
\emph{J.\ Math.\ Anal.\ Appl.}\ \textbf{354} (2009), 507--517.

\bibitem{Tao25}
T.~Tao,
``The maximal length of the Erd\H{o}s--Herzog--Piranian lemniscate in high degree,''
arXiv:2512.12455, December 2025.

\end{thebibliography}

\end{document}
